% EGA V, English translation by P. Blass, 1987 (cote 24): archive cote 24, site page 50 (PDF page 51).
% typescript. Draft compared: none.
% AI-transcribed by Claude Opus 5.5 (claude-opus-5-5, adaptive thinking, effort high), Message Batches API, batch ega24b (msgbatch_01RGJJeHdR96U16ZWFuUjtDx), 2026-11-13; not yet checked by eye.
\scanpage{24/50}
\mspage{50}

$U \neq 0$ Ysing / U $\Rightarrow$ U is an open immersion a fortiori $\underline{Y}$sing / U does \margin{?}
not contain the points of Ysupsing.  If X is proper over K, then
g : $\underline{Y}$sing $\to$ T is surjective thus $\underline{Y}$sing / U $\Rightarrow$ U is proper over K so
that g : $\underline{Y}$sing $\to$ T is surjective therefore $\underline{Y}$sing / U $\Rightarrow$ U is an
isomorphism and U is the biggest open set of T having the latter \margin{?}
property.

First of all since g is dominating and generically étale it
is generically étale at least one non empty open subset V of T
such that $\underline{Y}$sing $|$ V $\Rightarrow$ V is étale and surjective which implies that
V is smooth over K.  If then $\xi \in$ V(K) and if y is a point of $\underline{Y}$sing
(K) over $\xi$ then with the notations of 8.7 c) the space H$^1$ is
nothing else but the tangent space to T and $\xi$, and as we observed
here this implies that the point x of X(K), the projection of y is
determined as the orthogonal point to H$^1$ thus it is uniquely
determined thus since Ysing $\langle$ X x $\underline{P}$v, y is uniquely determined.

This proves already that g is birational (being generically
étale and generically radicial).  On the other hand the morphism $\psi$
(whose definition in its form is evident) which associates to
every $\xi \in$ U(K) the unique point x  $= \psi(\xi) \in \underline{P}$ orthogonal to the
tangent space to U at $\xi$, coinciding over V with the composition V
$\to \underline{Y}$sing $|$ V $\to$ X, where the second arrow is the projection;
therefore setting h = $(\psi$, id):U $\to \underline{P}$ x $\underline{T}$ (illegible)*
g$_1$ = g $|$ g$^{-1}$(U) : g$^{-1}$(U) $\to$ U the composition hg$_1$ : g$^{-1}$(U) $\to \underline{P}$ x T
is nothing else but the canonical inclusion, this being so for its
restriction to g$^{-1}$(V) $\xrightarrow{\sim}$ V.  It follows that h
factors through the scheme theoretic closure Y$^1$ of $\overline{\text{Y}^1}$ in $\underline{P}$ x T

\margin{* Ask AG}
